\section{Ruta futura: ¿convergerán Linear y Sparse?}

Esta sección examina en conjunto las rutas anteriores. En el programa se menciona una solución ideal para el futuro: combinar Linear Attention y Sparse Attention. Intuitivamente, Linear/Hybrid se encarga de reducir la KV cache y el costo recurrente de la mayoría de las capas, mientras que Sparse se encarga de seleccionar unos pocos token clave cuando se necesita una recuperación global. En teoría, si la selección sparse es precisa, puede sustituir una parte de las capas de Full Attention.

\subsection{Por qué resulta atractiva una ruta combinada}

Esta subsección explica por qué es probable que en el futuro no gane una sola ruta. Linear, Sparse y Full resuelven cada una un cuello de botella distinto, y los sistemas de ingeniería reales suelen combinarlas.

Linear Attention destaca por ahorrar estado y reducir el costo de la mayoría de los token, pero puede sacrificar algunas interacciones globales complejas; Sparse Attention conserva la capacidad de «mirar las posiciones importantes del historial», pero necesita un indexer suficientemente confiable. La ruta combinada busca que ambas se complementen: la mayor parte del tiempo se usa cheap recurrence y, cuando hace falta, sparse retrieval completa la información de largo alcance.

\begin{knowledgebox}{Tres posibles rutas combinadas}
\begin{tabularx}{\textwidth}{p{0.24\textwidth}p{0.34\textwidth}X}
\toprule
Forma de combinación & Mecanismo & Riesgo principal \\
\midrule
Linear + Full & Mayoría de capas lineales, pocas capas globales & Las capas globales todavía pueden limitar la context window. \\
Sparse + Full & Mayoría de capas dispersas, pocas capas globales & La KV cache sigue siendo grande y el indexer debe ser confiable. \\
Linear + Sparse & Estado lineal más recuperación dispersa & Mayor dificultad para la estabilidad del entrenamiento y mayor complejidad de implementación. \\
\bottomrule
\end{tabularx}
\end{knowledgebox}

\subsection{Qué significa «seleccionar con precisión»}

Lo esencial de Sparse Attention no es mirar menos token, sino mirar los correctos aun mirando menos. En el razonamiento de múltiples saltos, la recuperación en documentos largos, las dependencias de código y las trayectorias de agent, la información clave puede aparecer en posiciones muy lejanas. Si el indexer selecciona mal, se ahorra cómputo, pero el modelo responde mal; si selecciona demasiado, se pierde la ventaja de la dispersión.

\begin{warningbox}{Modos de falla de Sparse}
La falla más peligrosa de la atención dispersa no es la lentitud, sino omitir en silencio token clave. En las tareas de contexto largo, omitir una condición temprana o una restricción lejana puede hacer que todo el razonamiento posterior sea erróneo.
\end{warningbox}

\subsection{Resumen de la sección}

Es posible que el futuro de Attention no consista en elegir una sola entre Linear, Sparse y Full, sino en combinarlas. La verdadera dificultad es lograr que la ruta combinada conserve al mismo tiempo capacidad expresiva, bajo costo, adecuación al hardware y estabilidad del entrenamiento.
